
%% This is about discussing the results and answering the hypothesis. In very brief: The hypothesis whas that for tributary affected, step length is more important etc. Of course the discussion kind of depends on the results I find. One thing I can say for certain howerer is that one critique will be the small sample size. And also the lack of large non tributary affected patches.



% the method is not perfectly ideal and is acutally more of a screening method


% the buffer pix 


% the classification and selection of cwp is rather subjective. => So there is a lot of potential for confirmation bias. We assumed the hypothesis to be correct so this might has clouded my judgement. Next time it would defenetly be better to derive an algorithmic approach. But it should also be noted here that the design of a algorithmc set of rules of course is also succeptible to confirmation bias.


\section{Difference in cwp polygons between old and new function}

% the cwp patches produces by the old and the speed optimized function are generally comparable in their position and extent, but are not completely equivalent on the pixel level. Generally the optained jackard similarities are high. This is not optimal, however with the respect to the research quesitons it can be tollerated. The research question generally asks for the effect of the step length parameter on tributary caused or non tributary caused cwp. It is therefore not completely bound to the original algorithm, but more to the concept of cwp. Aslong as the optimized algorithm does a decent job at detecting and mapping out cwp, this is totaly fine and generally n